\section{Balancing the exterior tail of divisor feedback (NS-79)}
\label{sec:ns79-balanced-feedback}

The sampling norm isolates the exterior coefficient $\eta$. We therefore
test whether removing its change repairs the failed NS-75 step. This is
an explicitly modified correction, not a reinterpretation of that test.
Let $a_n^{(p)}=-r_N(n)\log^p(2N/n)$ be the three original new-block
physical directions, with the constant convention for $p=0$. Define
\begin{equation}
 \widetilde a_n^{(p)}=a_n^{(p)}\quad(N<n<2N),\qquad
 \widetilde a_{2N}^{(p)}=-2N\sum_{n=N+1}^{2N-1}\frac{a_n^{(p)}}n.
 \label{eq:ns79-balance}
\end{equation}

\begin{proposition}[Exact tail balance and its scope]
\label{prop:ns79-balance}
Each raw correction $f_p=\sum_{n=N+1}^{2N}\widetilde a_n^{(p)}\sigma_n$
vanishes for $x>1/(N+1)$. Thus the held-old step preserves $\eta$.
For $p=0$ it cancels the new derivative jumps for $N<m<2N$, leaving
the jump at $m=2N$ unconstrained. It is not an exact cancellation of
the entire next block.

The original three-direction span augmented by the physical endpoint
atom $\sigma_{2N}$ equals the balanced three-direction span augmented
by that atom. This identity persists after the complete old-space
projection. No assertion of norm descent follows from these identities.
\end{proposition}
\begin{proof}
The weighted coefficient sum in~\eqref{eq:ns79-balance} is zero.
For $x>1/(N+1)$ every new atom equals $1/(nx)$, proving vanishing
and tail preservation. In the open block each proper divisor is old,
so the NS-75 jump argument applies at every index except the modified
last one. More explicitly, for $N<t<2N$ the held-old updated residual
has one formula $-\eta t+D_N\log t-L_N$, with $D_N,L_N$ from NS-78;
no new jumps occur inside this open interval.

The difference between each original and balanced physical vector is
a multiple of the same endpoint coordinate. Hence augmenting either
span by that coordinate gives the same space; applying any linear
old-space projection preserves this equality. In adjacent-difference
coordinates the endpoint atom has vector $v_n=n/(2N)$, modulo its
old boundary term. The transformation is retained exactly in the
finite calculation.
\end{proof}

Old-space reoptimization can change $\eta$ again. The projected gains
below therefore do not preserve the raw tail balance. Conversely,
repeating only raw tail-preserving steps with a fixed nonzero initial
$\eta$ can never reduce squared error below $\eta^2$, because the
exterior tail remains unchanged. This observation does not exclude
initialization with $\eta=0$ or algorithms that reoptimize old coefficients.

\paragraph{Certified finite outcome.}
At 256/384 bits, with a doubled complete-kernel cutoff at $N=256$,
the fractions of full projected gain have the following strict enclosures.
\begin{center}
\begin{tabular}{rccc}
\hline
$N$&Balanced three-span&Endpoint-augmented four-span&Diagonal curvature\\
\hline
16&(.5960,.5961)&(.8542,.8543)&(.9308,.9310)\\
32&(.6580,.6581)&(.8240,.8241)&(.9040,.9041)\\
64&(.8978,.8979)&(.9192,.9193)&(.9783,.9784)\\
128&(.8809,.8810)&(.8810,.8811)&(.9852,.9854)\\
256&(.5996,.5997)&(.7118,.7119)&(.9763,.9764)\\
\hline
\end{tabular}
\end{center}
Thus removing the tail change improves the $N=256$ three-span fraction
from $56.02\%$ to $59.97\%$, and the full four-span reaches $71.18\%$;
both remain below curvature's $97.63\%$. The raw untapered unit step
at $N=256$ still increases squared error by a factor in $(53.07,53.09)$,
although this is smaller than NS-75's factor above $1495$.
Thirteen of the fifteen raw balanced unit steps increase error. The two
exceptions are the linear and quadratic tapers at $N=128$, whose error
ratios lie in $(.9643,.9645)$ and $(.9890,.9893)$. This is neither a
uniform-descent theorem nor an exclusion of all balanced feedback.
Sixteen exact Maxima checks cover endpoint compensation, coordinate
conversion and the complete exterior floor. The four-dimensional span
keeps all cross terms; no mixed term is discarded to claim positivity.
